% Internal pre-results manuscript.  This file is intentionally separate from
% main.tex, which preserves the historical allocation/router draft.
\documentclass[sigconf,anonymous]{aamas}
\usepackage{amsmath}
\usepackage{booktabs}
\usepackage{url}
% Copyright and conference block copied from the official AAMAS 2027 sample.
\makeatletter
\gdef\@copyrightpermission{
  \begin{minipage}{0.2\columnwidth}
   \href{https://creativecommons.org/licenses/by/4.0/}{\includegraphics[width=0.90\textwidth]{by}}
  \end{minipage}\hfill
  \begin{minipage}{0.8\columnwidth}
   \href{https://creativecommons.org/licenses/by/4.0/}{This work is licensed under a Creative Commons Attribution International 4.0 License.}
  \end{minipage}
  \vspace{5pt}
}
\makeatother
\setcopyright{ifaamas}
\acmConference[AAMAS 2027]{Proc.\@ of the 26th International Conference
on Autonomous Agents and Multiagent Systems (AAMAS 2027)}{3--7 May 2027}{Hanoi, Vietnam}{M.~Baldoni, F.~Fang, W.~Yeoh, N.~Yorke-Smith (eds.)}
\copyrightyear{2027}
\acmYear{2027}
\acmDOI{}
\acmISBN{}
\acmSubmissionID{INTERNAL REVISION}
\submissionType{Research Paper Track}

\acmSubmissionID{INTERNAL PRE-RESULTS -- NOT FOR SUBMISSION}

\title[Know Who You Are]{Know Who You Are: Earning Roles from Situated Peer Judgments}
\author{Anonymous Authors}
\affiliation{\institution{Internal research draft}}
\email{omitted}
\ccsdesc[500]{Computing methodologies~Multi-agent systems}
\ccsdesc[300]{Computing methodologies~Cooperation and coordination}
\keywords{multi-agent systems, self-evolving roles, peer judgment, online adaptation}

\begin{abstract}
Multi-agent systems often assign work using fixed roles, global reputation, or
terminal rewards. These signals can rank agents, but they cannot answer a sharper
question: did this particular delivery help the particular agent that depended on
it, and should that evidence change who receives the next task? We introduce
Earning Roles, a self-evolving responsibility mechanism that turns a recipient's
use, revision, or rejection of a versioned delivery into evidence tied to the
producer and task context. A responsibility check separates producer quality from
recipient integration and downstream outcomes; later independent outcomes can
revise the evidence before it informs another assignment. This gives each
collaboration a local, attributable learning signal while allowing useful
experience to transfer across recipients and tasks. We evaluate the mechanism on
multi-agent tasks with genuine dependencies and delayed feedback, comparing it
with static assignment, raw acceptance, contextual trust or bandit selection,
terminal-only learning, and pooled history under matched information and cost. We
measure whether it improves held-out task quality and total collaboration cost,
while retaining prior competence under drift. These tests distinguish role
learning from merely selecting a globally strong agent.
\end{abstract}

\begin{document}
\maketitle
\noindent\textbf{Internal pre-results manuscript.} The scientific release gate is closed.
All future result cells remain empty until an independently frozen experiment supplies
the corresponding raw evidence.

\section{Introduction}

An agent can be useful in one collaboration and unsuitable for another. The reason is
not only a latent capability: the recipient, dependency, artifact contract, and repair
responsibility determine whether a delivery helps. Existing multi-agent systems commonly
use fixed roles, a global reputation, or a central selector. Such mechanisms can rank
agents, but they do not establish whether a particular recipient used a particular
delivery, whether the delivery caused the later action, or whether a reported repair was
actually the producer's responsibility.

This paper asks one question: \emph{when can a real recipient's situated judgment become
reliable evidence for a producer's future responsibility?} The question requires a
complete causal chain:

\begin{center}
delivery $\rightarrow$ recipient judgment and use $\rightarrow$
attributable role evidence $\rightarrow$ pre-execution assignment $\rightarrow$
later quality and cost.
\end{center}

A label attached to a message, a final task score, or a change in a role posterior is not
enough. A useful mechanism must preserve the recipient's local information while keeping
the producer, recipient, and final evaluator accountable for different parts of the
outcome. This concern is related to reputation and partner-selection work
\cite{ren2026reputation,leung2025partner,russell2026partner}, but the central object here
is the attributable delivery and the later responsibility opportunity.

We contribute a protocol and a candidate update contract rather than an unsupported
claim of a new optimizer. The protocol fixes selection before execution, selected-only
delivery, delayed feedback, version binding, responsibility-aware labels, and complete
cost accounting. The candidate update, called responsibility-aware role evidence (RARE),
updates public evidence only when a recipient action can be attributed to the producer's
delivery. Raw acceptance, contextual trust/bandit selection, terminal-only feedback, and
a pooled controller provide strong alternative explanations. The paper's claim is
conditional: RARE is useful only if it changes a later legal assignment and improves an
unseen task stream beyond these matched alternatives.

The current evidence is deliberately limited. Two complete real-API development chains
show that selection, recipient judgment, source snapshots, consumer execution, and a later
assignment can be logged together. The same traces show that a four-check consumer score
missed producer priority defects, that normal recipient integration was mistaken for an
upstream repair, and that the feedback left the selection probability unchanged. These
observations determine the next experiment; they are not a positive method result.

\section{Problem and Information Contract}

At decision $t$, $x_t$ is the task context and $C_t$ is the legal set of candidate
producers. The selector samples $a_t$ before any candidate executes:
\begin{equation}
 a_t \sim \pi(\cdot\mid x_t,C_t,S_t), \qquad a_t\in C_t,
 \label{eq:select}
\end{equation}
where $S_t$ is the public role-evidence state available at that time. Only $a_t$ runs.
The resulting delivery is $o_t$. The real owner or downstream recipient is $j_t$ and
must differ from $a_t$ for a peer-judgment event. The recipient's observation arrives
after a delay $d_t$ and is bound to the decision-time candidate version.

The event record is intentionally small:
\begin{equation}
 e_t=(x_t,C_t,a_t,o_t,j_t,z_t,u_t,q_t,d_t,c_t),
 \label{eq:event}
\end{equation}
where $z_t$ is the recipient judgment, $u_t$ is the recipient action, $q_t$ is an
independent producer-contract or final-quality observation, and $c_t$ is the complete
cost record. The judgment and action are distinct. We use the labels
\texttt{accept}, \texttt{integrate}, \texttt{upstream\_fix}, \texttt{reject}, and
\texttt{unknown}; an unknown responsibility is missing information, not a negative label.
The record also stores the behavior probability $p_t=\pi(a_t\mid x_t,C_t,S_t)$,
artifact digest, candidate version, feature-schema version, and feedback watermark.

The online learner may read only events whose feedback has arrived. A hidden evaluator,
future task score, or unexecuted candidate outcome cannot affect the current choice,
retry, stopping rule, or update. Independent scores may be used after the stream for
evaluation; they are not online labels unless the benchmark explicitly makes the check
public and gives it to every condition.

\section{Candidate Mechanism: RARE}

The candidate mechanism maintains public evidence for producer $v$ in context $r$ as
\begin{equation}
 R(v,r)=(n(v,r),q(v,r),u(v,r)).
 \label{eq:state}
\end{equation}
Here $n$ is the amount of eligible evidence, $q$ is recipient evidence that can be
attributed to the producer delivery, and $u$ is evidence that is unresolved or rejected
by an independent responsibility check. An event increments $q$ only when the recipient
actually inspected the delivery, performed a traceable action because of it, and the
action is outside the recipient's ordinary integration responsibility. Otherwise it
increments $u$ or remains unknown. This rule prevents a consumer's own missing file or
normal integration work from becoming a producer penalty.

Given a cached representation $h(x,v,r)$, the candidate score is
\begin{equation}
 s(v\mid x,r)=w^\top h(x,v,r)+\beta
 \frac{q(v,r)-u(v,r)}{n(v,r)+\alpha}.
 \label{eq:score}
\end{equation}
The selector uses a predeclared exploration floor and records its propensity. When a
feedback event arrives, only the matching producer, context, and version receive one
increment. Delayed events are combined using event-level sufficient statistics; if they
cannot be combined without order dependence, the implementation replays the sealed event
window and records the replay watermark. The old-task holdout rejects a write whose
predeclared risk budget is exceeded.

RARE is a candidate mechanism, not evidence that this formula is optimal. In particular,
the same event stream is given to a raw-acceptance rule and to a contextual trust or
bandit baseline. If those baselines reproduce the future selections and utility, the
paper will report that the responsibility-aware mechanism added no identifiable value.
RLS, online logistic/SGD, and periodic refitting are training comparators; a small
incremental state alone is not a novelty claim.

\section{Benchmark and Evaluation Design}

The working benchmark is a TeamBench-derived PeerRoleBench-TB candidate. Its qualification
requires two structurally different task roots, an independently checked producer
contract, a genuine recipient dependency, and a root-level development/confirmation split.
The queue/priority DIST1 root is retained as a protocol diagnostic. The PIPE3 streaming
root remains a conditional second candidate because a parent process can separately check
producer output, recipient transformation, and downstream adoption. The current PIPE3
preflight is only a static contract/ledger fixture and listed-access canary; its task text
also exposes the intended serialization bugs and therefore cannot support discovery or
responsibility-attribution claims. A local adapter now constructs neutral task text and
redacted public source payloads, but its runtime dispatch is unverified. The benchmark is not locked until task text, runtime
payload/scorer isolation, and root-level split all qualify; renamed seeds are not independent
task roots.

The first comparison table is fixed before confirmation:
\begin{table}[t]
\centering
\small
\begin{tabular}{p{0.28\columnwidth}p{0.58\columnwidth}}
\toprule
Condition & Information and purpose\\
\midrule
Random/local uniform & Candidate set only; selection and exploration lower bound.\\
No update & Initial prior and current task; tests whether learning is needed.\\
Raw acceptance & Recipient accept/reject without responsibility separation.\\
Contextual trust/bandit & Same context, judgment, propensity, and budget; strongest simple explanation.\\
Terminal-only & Independent final outcome without situated recipient evidence.\\
RARE & Responsibility-aware recipient evidence; candidate mechanism.\\
Pooled controller & All permitted public history under the same total resource budget.\\
\bottomrule
\end{tabular}
\caption{Predeclared comparison boundary. Extra close methods are added only when their
implementation is faithful and their information and cost can be matched.}
\label{tab:baselines}
\end{table}

We ask four research questions. RQ1 tests whether a recipient judgment predicts independent
delivery and later-use outcomes beyond raw acceptance and terminal-only feedback. RQ2 tests
whether public evidence changes a future assignment made before execution by a recipient
that did not directly judge the candidate. RQ3 measures future prequential utility, full
cost, update latency, state size, and old-task forgetting under selected-only delayed
feedback. RQ4 tests responsibility ambiguity, recipient errors, version replacement, and
task drift.

The primary unit of uncertainty is an independently reset stream. Episodes within a stream
are dependent and are not treated as independent samples. The confirmation manifest fixes
task roots, seeds, prompts, models, candidate menus, exploration, token and call budgets,
failure handling, scorer versions, and the minimum useful effect before any confirmation
output is observed. All real calls, tool actions, retries, UNKNOWNs, and wall-clock costs
are streamed to auditable JSONL.

\section{Current Evidence and Limitations}

The completed v3 development run contains two limited real chains and six complete API
responses. Both recipients declared a repair, both changed only their own integration file,
and the recorded assignment was consumed by the next choice. The update mapping used the
same $0.5$ value as the prior, so the probability remained $[0.5,0.5]$ and the selected
peer did not change. The consumer score passed four checks but a separate diagnostic found
that both producer priority artifacts violated the public constructor or import contract.
These are protocol and scorer findings, not a role-learning measurement. The peers also
shared a fresh model and prompt with no persistent private experience, so their identity
labels cannot support a claim of learned specialization.

The consequence is a narrow current conclusion: the selected-only execution and audit
plumbing is usable after responsibility and scorer repairs, while the mechanism's signal,
peer differentiation, and efficacy remain unestablished. No A800 training run is justified
until a qualified task root produces an identifiable signal and a matched baseline contract.

\section{Discussion and Falsification}

The paper will withdraw the role-formation claim if a future assignment never changes, if
the assignment is made only by the directly judging owner, if a same-information trust or
bandit explains all gains, or if gains disappear at matched total cost. It will report a
protocol or negative result if recipient judgments are not independently informative. It
will not reinterpret a posterior change, a higher role entropy, or an isolated terminal
success as evidence of useful role formation.

This boundary matters for related systems that learn reputations or partner choices
\cite{ren2026reputation,leung2025partner,yang2025agentnet}. The proposed question is not
that agents can judge one another. It is whether a legally attributable, recipient-anchored
judgment changes a later responsibility opportunity and improves a full collaboration
objective under the same information and budget.

\section{Conclusion}

The current manuscript is a pre-results contract for a single empirical question. It
defines the observation boundary, a responsibility-aware candidate update, a benchmark
qualification path, and the strong explanations that must be defeated. The existing real
traces justify repairing measurement and identity state; they do not justify claiming a
working online role learner. The next scientific step is therefore a qualified second task
root and a small, frozen, matched comparison—not a larger model or an unbounded training
run.

\bibliographystyle{ACM-Reference-Format}
\bibliography{references}
\end{document}
